\documentclass{beamer}

\usepackage{../../../latex_style/beamerthemeExecushares}
\usepackage{../../../latex_style/notations}

\title{Session 7: Spectral Theorem, PCA and SVD}
\subtitle{Optimization and Computational Linear Algebra for Data Science}
\author{Léo Miolane}
\date{}

\setcounter{showSlideNumbers}{1}

\begin{document}
\setcounter{showProgressBar}{0}
\setcounter{showSlideNumbers}{0}

\frame{\titlepage}

\begin{frame}
	\frametitle{Contents}
	\begin{enumerate}
		\item The Spectral Theorem
		\item Principal Component Analysis
		\item Singular Value Decomposition
	\end{enumerate}
\end{frame}


\setcounter{framenumber}{0}
\setcounter{showSlideNumbers}{1}

\section{The Spectral Theorem}

\begin{frame}[t]{The spectral theorem}
	\grid

	\vspace{-0.3cm}
	\begin{block}{Theorem}
		Let $A \in \R^{n \times n}$ be a \textbf{symmetric} matrix. Then there is a orthonormal basis of $\R^n$ composed of eigenvectors of $A$.
	\end{block}

	\vspace{1.3cm}
	\begin{block}{Theorem (Matrix formulation)}
		Let $A \in \R^{n \times n}$ be a \textbf{symmetric} matrix. Then there exists an orthogonal matrix $P$ and a diagonal matrix $D$ of sizes $n \times n$ such that
		$$
		A = P D P^{\sT}.
		$$
	\end{block}
\end{frame}

\begin{frame}[t]{Geometric interpretation}
	\grid

\end{frame}

\begin{frame}[t]{A key result}
	\grid

	\vspace{-0.4cm}
	\begin{block}{Proposition}
		Let $A$ be a $n \times n$ symmetric matrix and let $\lambda_1 \geq \cdots \geq \lambda_n$ be its $n$ eigenvalues and $v_1, \dots, v_n$ be an associated orthonormal family of eigenvectors. Then 
		$$
		\lambda_1 = \max_{\|v\| = 1} v^{\sT} A v
		\qquad \text{and} \qquad
		v_1 = \argmax_{\|v\| = 1} v^{\sT} A v\,.
		$$
		Moreover, for $k=2,\dots,n$:
		$$
		\lambda_{k} = \max_{\|v\| = 1, \, v \perp v_1, \dots, v_{k-1}} v^{\sT} A v\,,
		\quad \text{and} \quad
		v_{k} = \argmax_{\|v\| = 1, \, v \perp v_1, \dots, v_{k-1}} v^{\sT} A v.
		$$
	\end{block}
\end{frame}

\begin{frame}[t]{Proof}
	\grid

	\pause
\end{frame}

\section{Principal Component Analysis}
\begin{frame}[t]{Empirical mean and covariance}
	\grid

	We are given a dataset of $n$ points $a_1, \dots, a_n \in \R^d$
	\vspace*{0.4cm}
	\begin{columns}
		\hspace*{-0.6cm}
		\begin{column}{0.45\textwidth}
			\vspace*{-1.7cm}
			\begin{center}
				\underline{$d=1$}
			\end{center}

			\begin{itemize}
				\item \textbf{Mean}
					$$
					\mu = \frac{1}{n}\sum_{i=1}^n a_i \quad \in \R
					$$
				\item \textbf{Variance}
					$$
					\sigma^2 = \frac{1}{n} \sum_{i=1}^n (a_i - \mu)^2 \quad \in \R
					$$
			\end{itemize}
		\end{column}
		\vrule
		\begin{column}{0.5\textwidth}
			\vspace*{-0.5cm}
			\pause
			\begin{center}
				\underline{$d \geq 2$}
			\end{center}
			\begin{itemize}
				\item \textbf{Mean}
					$$
					\mu = \frac{1}{n}\sum_{i=1}^n a_i \quad {\color{ExecusharesRed} \in \R^d}
					$$
				\item \textbf{Covariance matrix}
					\begin{align*}
						S &= \frac{1}{n} \sum_{i=1}^n (a_i - \mu)(a_i - \mu)^{\sT} \quad {\color{ExecusharesRed} \in \R^{d \times d}}
						\\
						  &= \frac{1}{n} \sum_{i=1}^n a_i a_i ^{\sT} \qquad {\color{ExecusharesRed} \text{if  } \mu=0.}
					\end{align*}
			\end{itemize}
		\end{column}
	\end{columns}
\end{frame}
\begin{frame}[t]{PCA}
	\grid

	\begin{itemize}
		\item We are given a dataset of $n$ points $a_1, \dots, a_n \in \R^d$, where $d$ is <<large>>.
			\vspace{0.1cm}
		\item \textbf{Goal:} represent this dataset in lower dimension, i.e.\ find $\widetilde{a}_1, \dots, \widetilde{a}_n \in \R^k$ where $k \ll d$.
			\vspace{0.1cm}
		\item Assume that the dataset is centered: $\sum_{i=1}^n a_i = 0$.
			\vspace{0.1cm}
		\item Then, $S$ can be simply written as:
			$$
			S = \sum_{i=1}^n a_i a_i^{\sT}
			= A^{\sT} A.
			$$
			where $A$ is the $n\times d$ ``data matrix'':
			$$
			A =
			\begin{pmatrix}
				- \ a_1 \ - \\
				\vdots  \\
				- \ a_n \ -
			\end{pmatrix}
			.
			$$
	\end{itemize}
\end{frame}

\begin{frame}[t]{Direction of maximal variance}
	\grid
	\pause

\end{frame}
\begin{frame}[t]{2nd direction of maximal variance}
	\grid

\end{frame}

\begin{frame}[t]{Which value of $k$ should we take?}
	\grid

\end{frame}


\section{Singular Value Decomposition}


\begin{frame}[t]{Singular values/vectors}
	\grid


\end{frame}
\begin{frame}[t]{Singular Value decomposition}
	\grid

	\begin{block}{Theorem}
		Let $A \in \R^{n \times m}$. Then there exists two orthogonal matrices $U \in \R^{n \times n}$ and $V \in \R^{m \times m}$ and a matrix $\Sigma \in \R^{n \times m}$ such that $\Sigma_{1,1} \geq \Sigma_{2,2}  \geq \cdots \geq 0$ and $\Sigma_{i,j} = 0$ for $i\neq j$, that verify
		$$
		A = U \Sigma V^{\sT}.
		$$
	\end{block}

\end{frame}

\begin{frame}[t]{Geometric interpretation of \ $U\Sigma V^{\sT}$}
\begin{figure}[H]
	\begin{center}
	\includegraphics[width = 0.6\linewidth]{../figures/svd.pdf}
	\end{center}
\end{figure}

\end{frame}

\appendix
\backupbegin
\begin{frame}[t]
	\frametitle{Questions?}
	\grid

	\pause
\end{frame}
\backupend

\end{document}
